\documentclass[11pt]{article}

\usepackage{amsmath,amsthm,amssymb}
\usepackage[margin=1in]{geometry}

\newtheorem{theorem}{Theorem}
\newtheorem{proposition}[theorem]{Proposition}
\newtheorem{corollary}[theorem]{Corollary}
\newtheorem{lemma}[theorem]{Lemma}
\theoremstyle{definition}
\newtheorem{definition}[theorem]{Definition}
\newtheorem{example}[theorem]{Example}
\newtheorem{remark}[theorem]{Remark}
\newtheorem{conjecture}[theorem]{Conjecture}

\newcommand{\Hom}{\operatorname{Hom}}
\newcommand{\End}{\operatorname{End}}
\newcommand{\id}{\mathrm{id}}
\newcommand{\C}{\mathcal{C}}
\newcommand{\Z}{\mathbb{Z}}

\title{When does a directed container decompose?\\
The two-atoms conjecture and its freeness correction}
\author{MacBeth\\\small Kodamai}
\date{9 June 2026}

\begin{document}
\maketitle

\begin{abstract}
A directed container is the same thing as a small category, and a recurring
dream is that every small category should split cleanly into two pure
``atoms'': a poset skeleton (the shape data, with no loops) and the
endomorphism monoids sitting at each vertex (the pure loop data), glued
together by composition as a mutual action --- a Zappa--Sz\'ep product. I make
this precise and then refute the naive form. The crisp result: for fixed
objects $a,b$ the source-oriented unique factorization $f = t \circ e$ (with
$e \in \End(a)$ an endomorphism at the source and $t$ a transversal arrow)
holds for every $f \in \Hom(a,b)$ \emph{if and only if} $\Hom(a,b)$ is a
\emph{free} right $\End(a)$-set; dually for the target orientation and free
left $\End(b)$-sets. A quick necessary numerical test in the finite case is
that $\lvert \End(a)\rvert$ divides $\lvert \Hom(a,b)\rvert$. The smallest
counterexample has five morphisms and fails this divisibility, so the
Zappa--Sz\'ep decomposition does not exist in general. The correct general
statement is a bimodule/profunctor assembly: each $\Hom(a,b)$ is an
$(\End(b),\End(a))$-biset, and $\C$ is the collage glueing of its vertex
endomorphism monoids along these connecting bisets. Zappa--Sz\'ep is exactly
the free, one-sided special case. This turns SEED open-question Q2 (the
distributive-law decision procedure) into a finitely checkable freeness
criterion.
\end{abstract}

\section{Setup: directed containers as small categories}

I take for granted the equivalence between directed containers and small
categories, and translate the dictionary once so that everything below is
phrased categorically.

\begin{definition}[Dictionary]
A small category $\C$ corresponds to a directed container as follows: the
\emph{shapes} are the objects; the positions $P(s)$ over a shape $s$ are the
morphisms out of $s$; the root position $o$ is the identity $\id_s$; the
operation $s \downarrow p$ takes the codomain of $p$; and the position monoid
operation $p \oplus q$ is composition \emph{with order reversed}, $p \oplus q =
q \circ p$. For objects $a,b$ I write $\Hom(a,b)$ for the set of morphisms
$a \to b$, and
\[
  \End(a) := \Hom(a,a)
\]
for the \emph{endomorphism monoid} at $a$, a monoid under composition with unit
$\id_a$.
\end{definition}

The two natural ways an endomorphism monoid acts on a hom-set are the whole
content of what follows, so I state them carefully.

\begin{definition}[The canonical biset structure]
Fix objects $a,b$. The monoid $\End(a)$ acts on the \emph{right} of $\Hom(a,b)$
by precomposition,
\[
  f \cdot e := f \circ e \qquad (f \in \Hom(a,b),\; e \in \End(a)),
\]
and the monoid $\End(b)$ acts on the \emph{left} by postcomposition,
\[
  e' \cdot f := e' \circ f \qquad (e' \in \End(b),\; f \in \Hom(a,b)).
\]
These two actions commute by associativity, $(e' \circ f) \circ e = e' \circ
(f \circ e)$, so $\Hom(a,b)$ is canonically an
$(\End(b),\End(a))$-\emph{biset}. The right action is unital and associative
because $\End(a)$ is a monoid; likewise the left.
\end{definition}

\section{The two atoms and the naive conjecture}

There are two ``pure'' kinds of small category, and the dream is that every
category is built from just these.

\begin{definition}[The two atoms]
\leavevmode
\begin{enumerate}
\item A \emph{poset} (more precisely a thin category): at most one arrow between
any ordered pair of objects, so in particular every endomorphism monoid is
trivial, $\End(a) = \{\id_a\}$. This is the pure \emph{shape} atom --- all
skeleton, no loops.
\item A \emph{monoid}: a one-object category, where every arrow is a loop. This
is the pure \emph{loop} atom --- all endomorphism, no shape.
\end{enumerate}
\end{definition}

\begin{conjecture}[Two atoms; SEED Q2, naive form]\label{conj:naive}
Every small category $\C$ is a Zappa--Sz\'ep product of a poset skeleton and
its vertex endomorphism monoids, glued by composition as the mutual action.
Concretely (\emph{source orientation}): there is a sub-collection of
\emph{transversal} arrows, forming a thin subcategory, such that every $f \in
\Hom(a,b)$ factors \emph{uniquely} as
\[
  f = t \circ e, \qquad e \in \End(a),\; t \text{ transversal}.
\]
Dually (\emph{target orientation}) every $f$ factors uniquely as $f = e' \circ
t$ with $e' \in \End(b)$.
\end{conjecture}

The conjecture bundles three claims: (a) such a factorization exists; (b) it is
unique; (c) the transversal is canonically the poset/reachability skeleton.
I will show all three can fail, pin the exact condition for (a)+(b), and
exhibit a separate failure of (c).

\section{The freeness correction}

Recall that a right $M$-set $X$ (for a monoid $M$) is \emph{free on a basis}
$T \subseteq X$ if the action map
\[
  \mu : T \times M \longrightarrow X, \qquad \mu(t,m) = t \cdot m
\]
is a bijection; equivalently $X \cong \coprod_{t \in T} M$ as right $M$-sets,
each summand the regular action. $X$ is \emph{free} if it admits some basis. For
a \emph{group} $M$ this is the same as the orbit-theoretic notion (every
stabilizer trivial, with $T$ one representative per orbit), but for a general
monoid the action-map formulation is the correct definition. A
\emph{transversal} is precisely a basis: one representative whose translates
$t \cdot M$ cover $X$ disjointly and faithfully.

\begin{theorem}[Freeness criterion]\label{thm:main}
Fix objects $a,b$ of a small category $\C$. The source-oriented unique
factorization of Conjecture~\ref{conj:naive} ---
\emph{every $f \in \Hom(a,b)$ factors uniquely as $f = t \circ e$ with $e \in
\End(a)$ and $t$ ranging over a fixed transversal} --- exists if and only if
$\Hom(a,b)$ is a \emph{free} right $\End(a)$-set. In that case a transversal is
exactly a choice of one representative arrow per $\End(a)$-orbit. Dually, the
target-oriented factorization $f = e' \circ t$ with $e' \in \End(b)$ exists iff
$\Hom(a,b)$ is a free left $\End(b)$-set.
\end{theorem}

\begin{proof}
I prove the source statement; the target statement is the mirror image under
$f \mapsto f^{\mathrm{op}}$.

Unfold the factorization. A transversal $T \subseteq \Hom(a,b)$ together with
the right $\End(a)$-action gives the action map
\[
  \mu : T \times \End(a) \longrightarrow \Hom(a,b), \qquad
  \mu(t,e) = t \circ e = t \cdot e.
\]
The factorization claim ``every $f$ is $t \circ e$ for some transversal $t$ and
some $e \in \End(a)$'' is exactly surjectivity of $\mu$; ``uniquely'' is
injectivity of $\mu$. Hence: \emph{source-oriented unique factorization at
$(a,b)$ holds $\iff$ there exists a transversal $T$ for which $\mu$ is a
bijection.}

But ``$\mu$ is a bijection for some $T$'' is, verbatim, the definition of
$\Hom(a,b)$ being a free right $\End(a)$-set with basis $T$. So the two
statements are literally the same statement, and the theorem holds. I record the
two directions explicitly for the reader.

($\Leftarrow$) If $\Hom(a,b)$ is free with basis $T$, then $\mu$ is a bijection
by definition; reading off $f = \mu(t,e) = t \cdot e$ gives existence
(surjectivity) and uniqueness (injectivity).

($\Rightarrow$) Conversely, if some transversal $T$ makes $\mu$ a bijection,
then $\Hom(a,b) \cong \coprod_{t \in T}\End(a)$ as right $\End(a)$-sets via
$\mu$, which is exactly freeness with basis $T$.

The phrase ``a transversal is a choice of one representative per orbit'' is the
group intuition; for a general monoid the precise statement is that a basis $T$
indexes the summands of the free decomposition, and surjectivity of $\mu$ says
$T$ meets every right $\End(a)$-orbit while injectivity says each restricted
orbit map $e \mapsto t\cdot e$ is injective \emph{and} distinct basis elements
have disjoint orbits.
\end{proof}

\begin{remark}
The slickest phrasing: a transversal making $\mu$ bijective is the same datum as
a free-basis of $\Hom(a,b)$ as a right $\End(a)$-set, and a set has a basis as a
free $M$-set iff the action is already free. ``Existence of factorization'' is
the transversal meeting every orbit; ``uniqueness'' is freeness (no two
endomorphisms send a transversal arrow to the same morphism).
\end{remark}

\begin{corollary}[Numerical necessary test]\label{cor:div}
If $\C$ is finite and the source-oriented unique factorization holds at $a,b$,
then $\lvert \End(a)\rvert$ divides $\lvert \Hom(a,b)\rvert$. Dually
$\lvert \End(b)\rvert$ divides $\lvert \Hom(a,b)\rvert$ for the target
orientation.
\end{corollary}

\begin{proof}
A free right $\End(a)$-set is a disjoint union of orbits each of size $\lvert
\End(a)\rvert$, so $\lvert \Hom(a,b)\rvert = (\text{number of orbits}) \cdot
\lvert \End(a)\rvert$.
\end{proof}

\section{The smallest counterexample}

\begin{example}[Five morphisms, divisibility fails both ways]\label{ex:main}
Let $\C$ have two objects $a,b$. Put $\Hom(a,b) = \{s\}$ a single connecting
arrow, $\Hom(b,a) = \varnothing$, and
\[
  \End(a) = \End(b) = \Z/2 = \{1, g\}, \qquad g^2 = 1,
\]
with the loops fixing the connecting arrow on both sides:
\[
  g \circ s = s = s \circ g.
\]
There are five morphisms in all: $\id_a = 1_a$, $g_a$, $\id_b = 1_b$, $g_b$,
and $s$.
\end{example}

\begin{proposition}
Example~\ref{ex:main} is a genuine small category, and it violates the source-
\emph{and} target-oriented unique factorization simultaneously.
\end{proposition}

\begin{proof}
\emph{It is a category.} The only composites to define beyond the two copies of
$\Z/2$ are those involving $s$. Since $s$ is the unique arrow $a\to b$ and there
are no arrows $b \to a$, every legal composite with $s$ lands in $\{s\}$:
$g_b \circ s, \; s \circ g_a, \; g_b \circ s \circ g_a$ are all forced to be
$s$, which is exactly the prescription $g\circ s = s = s\circ g$ (the loops are
``central'' on $s$ and fix it). Associativity holds because each $\Z/2$ is
associative and any bracketing of a string containing $s$ evaluates to $s$
(every loop fixes $s$), independent of the bracketing; unit laws are immediate.

\emph{It fails unique factorization.} Source orientation: $\End(a) = \{1,g\}$
has order $2$ but $\lvert \Hom(a,b)\rvert = 1$, and $2 \nmid 1$, so by
Corollary~\ref{cor:div} source-UF fails. Directly: $s \circ 1_a = s = s \circ
g_a$, so the two would-be distinct factorizations $s = t\circ 1$ and $s = t
\circ g$ (with $t = s$ the only candidate transversal) collapse --- the right
$\Z/2$-action on the singleton $\{s\}$ is the trivial action, which is not free.
Target orientation is symmetric: $1_b \circ s = s = g_b \circ s$, the left
action is trivial, and again $2 \nmid 1$. So both orientations fail at once.
\end{proof}

\begin{remark}[The one-sided minimal case already breaks]
Even the asymmetric category with $\End(a) = \Z/2$, $\End(b)$ trivial, and a
single arrow $a \to b$ fixed by the source loop already fails source-UF: it is
the smallest non-poset extension of the walking arrow ($\bullet \to \bullet$),
and one nontrivial loop acting trivially on a singleton hom-set is enough to
destroy uniqueness.
\end{remark}

\section{Two further nuances}

\begin{remark}[Unique factorization does not give a poset skeleton]\label{rem:thick}
Even when unique factorization holds, the transversal is a genuine \emph{choice}
and need not be the reachability poset --- it can be \emph{thick}. Let $\End(a)$
be trivial, $\End(b) = \Z/3$, and $\Hom(a,b)$ a set of three arrows on which
$\Z/3$ acts \emph{freely} on the left (a single regular orbit). Since $\End(a)$
is trivial, source-UF holds trivially. But the source-transversal is the entire
$\Hom(a,b)$: three parallel arrows $a \to b$, manifestly \emph{not} thin. So
clause (c) of the naive conjecture --- ``skeleton $=$ poset reflection'' ---
fails independently of clauses (a),(b).
\end{remark}

\begin{remark}[Source and target loops are independent]
Source loops act on the right and target loops on the left, and these are
logically independent: in general one cannot absorb both into a single
one-sided factorization. The naive conjecture conflates the two orientations,
whereas $\Hom(a,b)$ genuinely carries \emph{two} actions and freeness on one
side says nothing about the other. (In Example~\ref{ex:main} both fail; in
Remark~\ref{rem:thick} the right side is trivially free while the left is free
but with a thick transversal.)
\end{remark}

\begin{remark}[Computational evidence]
A brute-force script \texttt{two\_atoms\_check.py} enumerates roughly $26$ small
categories and tests source-oriented unique factorization on each. Of these,
$18/26$ satisfy source-UF and $8/26$ fail; every one of the $8$ failures
involves a nontrivial loop acting non-freely on a hom-set too small to carry a
free orbit --- precisely the mechanism of Example~\ref{ex:main} and
Corollary~\ref{cor:div}. This matches Theorem~\ref{thm:main} exactly: failure
$\iff$ non-free connecting biset.
\end{remark}

\section{The correct general statement}

The naive picture is salvaged not as a unique factorization but as a
bimodule/profunctor assembly.

\begin{theorem}[Corrected decomposition]\label{thm:correct}
Let $\C$ be a small category with object set $\mathrm{ob}\,\C$.
\begin{enumerate}
\item \emph{(General assembly.)} $\C$ is reconstructed from the data
\[
  \big(\End(a)\big)_{a \in \mathrm{ob}\,\C}, \qquad
  \big(\Hom(a,b) \text{ as an } (\End(b),\End(a))\text{-biset}\big)_{a,b},
\]
together with the composition maps
$\Hom(b,c) \otimes_{\End(b)} \Hom(a,b) \to \Hom(a,c)$, as the
\emph{collage} (category of elements / profunctor glueing) of the vertex
endomorphism monoids along these connecting bisets. This is a
bimodule-style assembly, not a unique factorization, and it \emph{always}
exists.
\item \emph{(Zappa--Sz\'ep special case.)} $\C$ is a Zappa--Sz\'ep product of a
transversal skeleton and its vertex endomorphism monoids, in the sense of the
source-oriented Conjecture~\ref{conj:naive}, if and only if every connecting
biset $\Hom(a,b)$ is \emph{free} as a right $\End(a)$-set (the one-sided,
regular case). Dually for the target orientation and left freeness.
\end{enumerate}
\end{theorem}

\begin{proof}
(1) Composition in $\C$ restricted to $\End(a)$ recovers each vertex monoid, and
the biset structure of Definition~2 records pre/post-composition; the full
composition maps factor through the relative tensor product over $\End(b)$
because $(g \circ f)$ depends on $f,g$ only up to inserting/removing a $b$-loop,
which is exactly the coequalizer defining $\otimes_{\End(b)}$. Conversely this
data plus the composition maps reassembles $\C$: objects are the indices,
morphisms $a \to b$ are the biset elements, identities are the monoid units, and
the composition maps give associative composition (associativity is the
compatibility of the tensor glueing). This is the standard collage of the
endomorphism monoids regarded as a single-object categories linked by the
$\Hom$-profunctors.

(2) Apply Theorem~\ref{thm:main} at every pair $a,b$: source-oriented unique
factorization globally is exactly the statement that every connecting biset is a
free right $\End(\text{source})$-set, in which case picking a transversal per
orbit gives the (thin-or-thick) skeleton and the loops, with composition as the
mutual action --- a Zappa--Sz\'ep product. If some biset is non-free, the
factorization fails there by Theorem~\ref{thm:main}, so no global Zappa--Sz\'ep
decomposition exists.
\end{proof}

So the two-atoms dream is \emph{true exactly in the free, regular case} and
false in general; the universal truth is the weaker bimodule collage.

\section{Grant relevance}

This converts SEED open-question Q2 --- the distributive-law / Zappa--Sz\'ep
decision procedure for directed containers --- into a \emph{checkable
criterion}. A Zappa--Sz\'ep (equivalently, distributive-law) decomposition of a
directed container into a skeleton and its vertex monoids exists if and only if
each hom-set is free over the relevant vertex endomorphism monoid. Freeness is
finitely checkable (compute orbits, test that each orbit map is injective), and
the counting condition $\lvert \End(a)\rvert \mid \lvert \Hom(a,b)\rvert$ of
Corollary~\ref{cor:div} is a quick necessary pre-test that rejects most
non-decomposable cases instantly. For the grant this is a clean, implementable
algorithm with a sharp theorem behind it: the decision procedure is exactly an
orbit-freeness check on the connecting bisets, and when it fails, the corrected
decomposition theorem (Theorem~\ref{thm:correct}) still gives a canonical
profunctor-collage assembly to fall back on.

\end{document}
